% figures_arch.tex — architecture and pipeline diagrams (TikZ, no external assets).
% Topologies drawn here correspond to kinetic_ai/models/eqlm.py (EqLM, EqLMCore)
% and kinetic_ai/models/kinetic_lm.py (KineticLM conversion of a pretrained stack).

\begin{figure}[t]
\centering
\begin{tikzpicture}[
  font=\small,
  block/.style={draw, rounded corners=2pt, minimum width=2.1cm, minimum height=0.62cm, align=center},
  explicit/.style={block, fill=black!6},
  tied/.style={block, fill=blue!12, draw=blue!55},
  core/.style={block, fill=orange!16, draw=orange!65},
  io/.style={draw, rounded corners=6pt, minimum width=2.1cm, minimum height=0.55cm, fill=black!3},
  ar/.style={-{Latex[length=2mm]}, gray!70},
  lbl/.style={font=\scriptsize, align=center, text=black!65}
]

% ---- (a) explicit stack -------------------------------------------------
\begin{scope}[xshift=0cm]
  \node[io] (ain) at (0,0) {tokens};
  \foreach \i/\y in {1/0.95, 2/1.75, 3/2.55} {
    \node[explicit] (a\i) at (0,\y) {layer \i};
  }
  \node[lbl] at (0,3.25) {$\vdots$};
  \node[explicit] (aL) at (0,3.95) {layer $L$};
  \node[io] (aout) at (0,4.85) {logits};
  \draw[ar] (ain) -- (a1); \draw[ar] (a1) -- (a2); \draw[ar] (a2) -- (a3);
  \draw[ar] (a3) -- (2.15,2.55) -- (2.15,3.95) -- (aL);
  \draw[ar] (aL) -- (aout);
  \node[lbl, anchor=north] at (0,-0.55) {(a) explicit stack\\$L$ distinct blocks\\$O(L)$ params, $O(L)$ activations};
\end{scope}

% ---- (b) EqLM: single tied block solved to a fixed point ----------------
\begin{scope}[xshift=5.1cm]
  \node[io] (bin) at (0,0) {tokens};
  \node[tied, minimum height=1.5cm] (bblk) at (0,2.4) {$f_\theta(z,x)$\\\scriptsize one tied block};
  \node[io] (bout) at (0,4.85) {logits};
  \draw[ar] (bin) -- (bblk);
  \draw[ar] (bblk) -- (bout);
  \draw[ar, blue!65] (bblk.east) .. controls (2.3,2.9) and (2.3,1.9) .. (bblk.south east)
    node[lbl, text=blue!65, pos=0.5, right=1pt] {$z \leftarrow f_\theta(z,x)$};
  \node[lbl, anchor=north] at (0,-0.55) {(b) \EqLM: fixed point\\solved to $z^*=f_\theta(z^*,x)$\\$O(1)$ params, $O(1)$ activations};
\end{scope}

% ---- (c) KineticLM: explicit shell around a recursive core --------------
\begin{scope}[xshift=10.2cm]
  \node[io] (cin) at (0,0) {tokens};
  \node[explicit] (cpre1) at (0,0.85) {pre $1..n$};
  \node[core, minimum height=1.45cm] (ccore) at (0,2.4) {shared core\\\scriptsize applied $K$ times};
  \node[explicit] (cpost) at (0,3.95) {post $1..n$};
  \node[io] (cout) at (0,4.85) {logits};
  \draw[ar] (cin) -- (cpre1); \draw[ar] (cpre1) -- (ccore);
  \draw[ar] (ccore) -- (cpost); \draw[ar] (cpost) -- (cout);
  \draw[ar, orange!75] (ccore.east) .. controls (2.3,2.9) and (2.3,1.9) .. (ccore.south east)
    node[lbl, text=orange!75, pos=0.5, right=1pt] {$\times K$};
  \node[lbl, anchor=north] at (0,-0.55) {(c) \KineticLM: converted base\\outer layers kept explicit,\\middle tied and iterated};
\end{scope}

\end{tikzpicture}
\caption{Depth as an explicit stack, as a fixed point, and as a converted
hybrid. The explicit stack (a) spends parameters and activation memory linearly
in depth. \EqLM{} (b) replaces the stack with one weight-tied block iterated to
a fixed point, making activation memory constant in effective depth at the cost
of solving a nonlinear system each forward pass. \KineticLM{} (c) is the
conversion applied to a pretrained model: the first and last $n$ layers are
retained as initialized, and the middle layers collapse into a single shared
core applied $K$ times, which is the configuration the damage measurements of
Section~\ref{sec:results} select.}
\label{fig:topologies}
\end{figure}

\begin{figure}[t]
\centering
\begin{tikzpicture}[
  font=\small, node distance=0.55cm,
  stage/.style={draw, rounded corners=2pt, minimum width=3.1cm, minimum height=0.68cm, align=center, fill=black!4},
  gate/.style={draw, diamond, aspect=2.1, inner sep=1pt, align=center, fill=blue!8, draw=blue!55, font=\scriptsize},
  term/.style={draw, rounded corners=7pt, minimum width=2.4cm, minimum height=0.6cm, align=center, fill=black!3},
  ar/.style={-{Latex[length=2mm]}, gray!75},
  lbl/.style={font=\scriptsize, text=black!65}
]
  \node[term] (h) {hypothesis with\\numeric gate};
  \node[stage, right=1.1cm of h] (spec) {specification\\(arms, seeds, metric)};
  \node[stage, right=1.1cm of spec] (run) {execution\\(config hash + commit)};
  \node[gate, right=1.1cm of run] (rev) {adversarial\\review};
  \node[stage, below=1.15cm of rev] (amend) {claim rescoped\\or refuted};
  \node[gate, right=1.15cm of rev] (gate) {gate\\met?};
  \node[term, above right=0.35cm and 1.0cm of gate] (met) {\textsc{met}};
  \node[term, right=1.0cm of gate] (part) {\textsc{partial}};
  \node[term, below right=0.35cm and 1.0cm of gate] (miss) {\textsc{missed}};

  \draw[ar] (h) -- (spec);
  \draw[ar] (spec) -- (run);
  \draw[ar] (run) -- (rev);
  \draw[ar] (rev) -- (gate);
  \draw[ar] (rev) -- node[lbl, right] {refuted} (amend);
  \draw[ar] (amend) -| ($(run.south) + (-0.6,-0.75)$) -- ($(run.south) + (-0.6,0)$);
  \draw[ar] (gate) -- (met);
  \draw[ar] (gate) -- (part);
  \draw[ar] (gate) -- (miss);
  \node[lbl, below=0.1cm of amend, align=center] {a rescoped claim re-enters\\at execution, not at review};
\end{tikzpicture}
\caption{The closure protocol applied to every hypothesis reported here. A
quantitative gate is fixed before execution; each run records the SHA-256 of its
resolved configuration and the commit that produced it; an adversarial review
recomputes every number from raw results and may refute or rescope the claim
before it is scored. Outcomes are named rather than binary, so a documented miss
terminates a line of enquiry with the same standing as a pass.}
\label{fig:protocol}
\end{figure}
